\section{Class attributes}

\begin{frame}{What are class attributes?}
\pause
\begin{block}{Class attributes (definition)}
    Class attributes are \textbf{data attributes} or \textbf{methods}, that do not belong to instances of a class, but to the class itself.
\end{block}
\pause
\begin{itemize}
    \item Class data attributes
    \pause
    \item Class methods
\end{itemize}
\end{frame}

\begin{frame}[fragile]{An example...}
\pause
\begin{minted}{python}
class Person():
    ...

    # An oridinary method :o
    def greet(self):
        print("Hey!")
\end{minted}
\end{frame}

\begin{frame}[fragile]{Hmmmmm...}
\pause
\begin{minted}{python}
person1 = Person("Daan", "Wichmann", 19)
\end{minted}
\pause
\begin{minted}{python}
person2 = Person("Ella", "Kemperman", 19)
\end{minted}
\pause
\begin{minted}{python}
person1.greet()
>>> "Hey!"
\end{minted}
\pause
\begin{minted}{python}
person2.greet()
>>> "Hey!"
\end{minted}
\pause
\begin{itemize}
    \item Different objects, same output!
    \pause
    \item We can do this better!
\end{itemize}
\end{frame}

\begin{frame}[fragile]{Wow it's a classmethod!}
\pause
\begin{minted}{python}
class Person():
    ...

    @classmethod # Only the decorator is added!
    def greet(cls):  # Note the first parameter!
        print("Hey!")
\end{minted}
\end{frame}

\begin{frame}[fragile]{Same for data attributes!}
\pause
\begin{itemize}
    \item If the values are the same for every object, we can make it a class data attribute!
\end{itemize}
\begin{minted}{python}
class Person():
    greeting = "Hey! # Class data attribute

    @classmethod
    def greet(cls):
        print(cls.greeting)
\end{minted}
\end{frame}